% This section is written after Section~\ref{sec:results} has real numbers
% and docs/experiments_findings.md's error-analysis categories have been
% manually reviewed.

This section addresses three questions directly from the results in
Section~\ref{sec:results} and the error-analysis output in
\texttt{docs/experiments\_findings.md}.

\subsection{Does retrieval help overall, and is the effect reliable?}

At the default retrieval depth ($k=5$), the retrieval-augmented condition
outperformed the baseline on every metric: closed-question accuracy rose
from 43.0\% to 47.8\% (+4.8 points), overall exact match from 26.8\% to
29.5\% (+2.7 points), with smaller gains on BLEU-4 and ROUGE-L. However,
the paired bootstrap significance test on the closed-question delta
returned $p = 0.1848$ with a 95\% confidence interval of $[-0.020, 0.116]$,
and the open-question ROUGE-L delta returned $p = 0.9642$ with a
confidence interval of $[-0.032, 0.033]$ -- both spanning zero. With
$n = 451$ test examples (251 closed, 200 open), we cannot reject the null
hypothesis that either improvement arose by chance, though the
closed-question result is the more suggestive of the two. The honest
reading is that retrieval augmentation produced a consistent,
direction-positive effect across all reported metrics, but the test set
is too small, and the effect too modest, to confirm it is real rather
than sampling noise.

\subsection{Why does performance peak at $k=1$ rather than at higher retrieval depths?}

The top-$k$ sweep (Table~\ref{tab:topk-sweep}, Section~\ref{sec:results}) shows a
monotonic decline in closed-question accuracy as $k$ increases from 1 to 10:
accuracy falls from 55.0\% at $k=1$ to 46.6\% at $k=10$, while
open-question BLEU-4 and ROUGE-L do not show as clean a monotonic trend
(ROUGE-L is in fact highest at $k=10$, 0.1516). The clearest and most
interpretable effect in the sweep is on closed-question accuracy. A
plausible explanation is evidence dilution: mean top-1 retrieval
similarity was high (0.971), indicating the single most similar corpus
image is usually a strong match, but each additional retrieved caption
beyond the first introduces topically related but increasingly less
relevant text into the prompt. For a base (not instruction-tuned)
checkpoint such as \texttt{blip2-opt-2.7b}, a longer, noisier evidence
block may compete with the actual question for the model's limited
effective context rather than reinforcing it. An equally plausible
alternative explanation, not distinguishable from evidence dilution using
this experiment alone, is that the retriever itself is question-agnostic:
because it embeds only the query image (Section~\ref{sec:method}) and never
the question text, only the single most visually similar corpus entry is
reliably relevant to what is actually being asked, and captions beyond the
first are increasingly likely to be irrelevant to the specific question
rather than merely diluted -- a distinction future work with
question-aware retrieval (Section~\ref{sec:limitations}) could disentangle.

\subsection{What do the improved/regressed error categories show?}

Of the 451 test examples, retrieval changed the outcome for a similar
number of cases in each direction: 54 examples improved (baseline wrong,
RAG correct) versus 44 that regressed (baseline correct, RAG wrong), a
near-even split that is consistent with the modest, non-significant
aggregate delta reported above. The dominant category, however, was
\emph{both wrong} (260 of 451 examples, 58\%), with only 93 examples
(21\%) answered correctly by both conditions. This indicates that the
baseline VLM's zero-shot capability on radiology imagery, not the
presence or absence of retrieved evidence, is the primary bottleneck on
this benchmark: retrieval can only change the outcome on the minority of
examples where the underlying model was close to a correct answer in the
first place.
